% paper66-education-platform.tex
% JuGeo as a Teaching Tool: Interactive Proof Exploration for Students
% JuGeo Paper Series — Paper 66

\documentclass[11pt]{article}
% jugeo-common.tex — shared preamble for all Judgment Geometry papers
% Include via: % jugeo-common.tex — shared preamble for all Judgment Geometry papers
% Include via: % jugeo-common.tex — shared preamble for all Judgment Geometry papers
% Include via: \input{jugeo-common}

% ── Packages ────────────────────────────────────────────────────────────
\usepackage{amsmath,amssymb,amsthm,mathtools}
\usepackage{tikz-cd}
\usepackage{listings}
\usepackage{xcolor}
\usepackage{xspace}
\usepackage{booktabs}
\usepackage{multirow}
\usepackage{enumitem}
\usepackage[expansion=false]{microtype}
\usepackage{hyperref}
\usepackage{cleveref}
% \usepackage{thmtools}  % disabled: not available in this TeX installation
\usepackage{float}
\usepackage{caption}
\usepackage{subcaption}
\usepackage{tabularx}
\usepackage{stmaryrd}       % \llbracket, \rrbracket
\usepackage{bbm}            % \mathbbm{1}

% ── Page geometry ───────────────────────────────────────────────────────
\usepackage[margin=1in]{geometry}

% ── Theorem environments ────────────────────────────────────────────────
\theoremstyle{plain}
\newtheorem{theorem}{Theorem}[section]
\newtheorem{lemma}[theorem]{Lemma}
\newtheorem{proposition}[theorem]{Proposition}
\newtheorem{corollary}[theorem]{Corollary}

\theoremstyle{definition}
\newtheorem{definition}[theorem]{Definition}
\newtheorem{example}[theorem]{Example}
\newtheorem{construction}[theorem]{Construction}

\theoremstyle{remark}
\newtheorem{remark}[theorem]{Remark}
\newtheorem{notation}[theorem]{Notation}

% ── Python listings style ──────────────────────────────────────────────
\definecolor{jugeoBlue}{HTML}{2563EB}
\definecolor{jugeoGreen}{HTML}{059669}
\definecolor{jugeoRed}{HTML}{DC2626}
\definecolor{jugeoPurple}{HTML}{7C3AED}
\definecolor{jugeoGray}{HTML}{6B7280}
\definecolor{jugeoBg}{HTML}{F8FAFC}

\lstdefinestyle{jugeo-python}{
  language=Python,
  basicstyle=\ttfamily\small,
  keywordstyle=\color{jugeoBlue}\bfseries,
  stringstyle=\color{jugeoGreen},
  commentstyle=\color{jugeoGray}\itshape,
  numberstyle=\tiny\color{jugeoGray},
  backgroundcolor=\color{jugeoBg},
  numbers=left,
  numbersep=6pt,
  frame=single,
  framerule=0.4pt,
  rulecolor=\color{jugeoGray!40},
  breaklines=true,
  breakatwhitespace=true,
  showstringspaces=false,
  tabsize=4,
  xleftmargin=1.5em,
  framexleftmargin=1.5em,
  morekeywords={dataclass,frozen,field,Enum,IntEnum,auto,Any,
                Mapping,Sequence,Optional,match,case},
  literate={→}{{$\to$}}1 {←}{{$\leftarrow$}}1
           {≤}{{$\leq$}}1 {≥}{{$\geq$}}1 {≠}{{$\neq$}}1
           {∈}{{$\in$}}1 {∀}{{$\forall$}}1 {∃}{{$\exists$}}1
           {⊕}{{$\oplus$}}1 {⊖}{{$\ominus$}}1,
  escapeinside={(*@}{@*)},
}
\lstset{style=jugeo-python}

\lstdefinestyle{jugeo-output}{
  basicstyle=\ttfamily\small,
  backgroundcolor=\color{jugeoBg},
  frame=single,
  framerule=0.4pt,
  rulecolor=\color{jugeoGray!40},
  breaklines=true,
  showstringspaces=false,
  xleftmargin=1.5em,
  framexleftmargin=0.5em,
  numbers=none,
}

% ── JuGeo-specific macros ──────────────────────────────────────────────

% Judgment 8-tuple
\newcommand{\J}{\mathcal{J}}
\newcommand{\Jud}[8]{(#1,\,#2,\,#3,\,#4,\,#5,\,#6,\,#7,\,#8)}
\newcommand{\JudShort}[1]{\J_{#1}}

% Components
\newcommand{\coord}{\mathsf{c}}            % coordinate
\newcommand{\prop}{\varphi}                 % proposition
\newcommand{\carrier}{\mathcal{A}}          % carrier type
\newcommand{\evid}{\mathcal{E}}             % evidence bundle
\newcommand{\oblig}{\mathcal{O}}            % obligations
\newcommand{\obstr}{\mathcal{B}}            % obstructions
\newcommand{\trust}{\mathcal{T}}            % trust annotation
\newcommand{\prov}{\Pi}                     % provenance

% Trust algebra
\newcommand{\Talg}{\mathfrak{T}}            % trust ordered algebra
\newcommand{\Eadm}{\mathcal{E}_{\mathrm{adm}}}
\newcommand{\tleq}{\preceq}                 % trust ordering
\newcommand{\tjoin}{\oplus}                 % conservative join
\newcommand{\tatten}{\ominus}               % attenuation
\newcommand{\tpromote}{\uparrow_\pi}        % promotion
\newcommand{\tdemote}{\downarrow_\chi}       % demotion

% Trust tiers
\newcommand{\Tcontra}{\ensuremath{\mathsf{CONTRADICTED}}}
\newcommand{\Tunver}{\ensuremath{\mathsf{UNVERIFIED}}}
\newcommand{\Tcopilot}{\ensuremath{\mathsf{COPILOT}}}
\newcommand{\Toracle}{\ensuremath{\mathsf{ORACLE}}}
\newcommand{\Truntime}{\ensuremath{\mathsf{RUNTIME}}}
\newcommand{\Tsolver}{\ensuremath{\mathsf{SOLVER}}}
\newcommand{\Tproof}{\ensuremath{\mathsf{PROOF}}}

% Site and geometry
\newcommand{\Site}{\mathcal{S}}             % semantic site
\newcommand{\Cov}{\mathrm{Cov}}            % covering family
\newcommand{\Ob}{\mathrm{Ob}}              % objects
\newcommand{\Mor}{\mathrm{Mor}}            % morphisms
\newcommand{\res}{\mathrm{res}}            % restriction
\newcommand{\Cech}{\check{C}}              % Čech complex
\newcommand{\Hcoh}[1]{H^{#1}}             % cohomology group

% Descent
\newcommand{\descent}{\mathrm{desc}}
\newcommand{\glue}{\mathrm{glue}}
\newcommand{\overlap}[2]{#1 \cap #2}

% Semantic moves
\newcommand{\VERIFY}{\textsc{Verify}}
\newcommand{\CONSTRUCT}{\textsc{Construct}}
\newcommand{\NORMALIZE}{\textsc{Normalize}}
\newcommand{\GLUE}{\textsc{Glue}}
\newcommand{\RESTRICT}{\textsc{Restrict}}
\newcommand{\TRANSPORT}{\textsc{Transport}}
\newcommand{\REPAIR}{\textsc{Repair}}
\newcommand{\PROMOTE}{\textsc{Promote}}
\newcommand{\CHALLENGE}{\textsc{Challenge}}
\newcommand{\ESCALATE}{\textsc{Escalate}}

% Fragments
\newcommand{\QFLIA}{\texttt{QF\_LIA}}
\newcommand{\QFLRA}{\texttt{QF\_LRA}}
\newcommand{\QFBV}{\texttt{QF\_BV}}
\newcommand{\QFUF}{\texttt{QF\_UF}}

% Misc
\newcommand{\Python}{\textsc{Python}}
\newcommand{\JuGeo}{\textsc{JuGeo}}
\newcommand{\LEAN}{\textsc{Lean}\,4}
\newcommand{\Fstar}{F$^{\star}$}
\newcommand{\Coq}{\textsc{Coq}}
\newcommand{\Dafny}{\textsc{Dafny}}

% Common shortcuts
\newcommand{\ie}{i.e.\@\xspace}
\newcommand{\eg}{e.g.\@\xspace}
\newcommand{\cf}{cf.\@\xspace}
\newcommand{\wrt}{w.r.t.\@\xspace}

% Evaluation shortcuts
\newcommand{\numcases}{300}
\newcommand{\numspec}{100}
\newcommand{\numeq}{100}
\newcommand{\numbug}{100}
\newcommand{\medtime}{6.5\,ms}
\newcommand{\totaltime}{1.949\,s}


% ── Packages ────────────────────────────────────────────────────────────
\usepackage{amsmath,amssymb,amsthm,mathtools}
\usepackage{tikz-cd}
\usepackage{listings}
\usepackage{xcolor}
\usepackage{xspace}
\usepackage{booktabs}
\usepackage{multirow}
\usepackage{enumitem}
\usepackage[expansion=false]{microtype}
\usepackage{hyperref}
\usepackage{cleveref}
% \usepackage{thmtools}  % disabled: not available in this TeX installation
\usepackage{float}
\usepackage{caption}
\usepackage{subcaption}
\usepackage{tabularx}
\usepackage{stmaryrd}       % \llbracket, \rrbracket
\usepackage{bbm}            % \mathbbm{1}

% ── Page geometry ───────────────────────────────────────────────────────
\usepackage[margin=1in]{geometry}

% ── Theorem environments ────────────────────────────────────────────────
\theoremstyle{plain}
\newtheorem{theorem}{Theorem}[section]
\newtheorem{lemma}[theorem]{Lemma}
\newtheorem{proposition}[theorem]{Proposition}
\newtheorem{corollary}[theorem]{Corollary}

\theoremstyle{definition}
\newtheorem{definition}[theorem]{Definition}
\newtheorem{example}[theorem]{Example}
\newtheorem{construction}[theorem]{Construction}

\theoremstyle{remark}
\newtheorem{remark}[theorem]{Remark}
\newtheorem{notation}[theorem]{Notation}

% ── Python listings style ──────────────────────────────────────────────
\definecolor{jugeoBlue}{HTML}{2563EB}
\definecolor{jugeoGreen}{HTML}{059669}
\definecolor{jugeoRed}{HTML}{DC2626}
\definecolor{jugeoPurple}{HTML}{7C3AED}
\definecolor{jugeoGray}{HTML}{6B7280}
\definecolor{jugeoBg}{HTML}{F8FAFC}

\lstdefinestyle{jugeo-python}{
  language=Python,
  basicstyle=\ttfamily\small,
  keywordstyle=\color{jugeoBlue}\bfseries,
  stringstyle=\color{jugeoGreen},
  commentstyle=\color{jugeoGray}\itshape,
  numberstyle=\tiny\color{jugeoGray},
  backgroundcolor=\color{jugeoBg},
  numbers=left,
  numbersep=6pt,
  frame=single,
  framerule=0.4pt,
  rulecolor=\color{jugeoGray!40},
  breaklines=true,
  breakatwhitespace=true,
  showstringspaces=false,
  tabsize=4,
  xleftmargin=1.5em,
  framexleftmargin=1.5em,
  morekeywords={dataclass,frozen,field,Enum,IntEnum,auto,Any,
                Mapping,Sequence,Optional,match,case},
  literate={→}{{$\to$}}1 {←}{{$\leftarrow$}}1
           {≤}{{$\leq$}}1 {≥}{{$\geq$}}1 {≠}{{$\neq$}}1
           {∈}{{$\in$}}1 {∀}{{$\forall$}}1 {∃}{{$\exists$}}1
           {⊕}{{$\oplus$}}1 {⊖}{{$\ominus$}}1,
  escapeinside={(*@}{@*)},
}
\lstset{style=jugeo-python}

\lstdefinestyle{jugeo-output}{
  basicstyle=\ttfamily\small,
  backgroundcolor=\color{jugeoBg},
  frame=single,
  framerule=0.4pt,
  rulecolor=\color{jugeoGray!40},
  breaklines=true,
  showstringspaces=false,
  xleftmargin=1.5em,
  framexleftmargin=0.5em,
  numbers=none,
}

% ── JuGeo-specific macros ──────────────────────────────────────────────

% Judgment 8-tuple
\newcommand{\J}{\mathcal{J}}
\newcommand{\Jud}[8]{(#1,\,#2,\,#3,\,#4,\,#5,\,#6,\,#7,\,#8)}
\newcommand{\JudShort}[1]{\J_{#1}}

% Components
\newcommand{\coord}{\mathsf{c}}            % coordinate
\newcommand{\prop}{\varphi}                 % proposition
\newcommand{\carrier}{\mathcal{A}}          % carrier type
\newcommand{\evid}{\mathcal{E}}             % evidence bundle
\newcommand{\oblig}{\mathcal{O}}            % obligations
\newcommand{\obstr}{\mathcal{B}}            % obstructions
\newcommand{\trust}{\mathcal{T}}            % trust annotation
\newcommand{\prov}{\Pi}                     % provenance

% Trust algebra
\newcommand{\Talg}{\mathfrak{T}}            % trust ordered algebra
\newcommand{\Eadm}{\mathcal{E}_{\mathrm{adm}}}
\newcommand{\tleq}{\preceq}                 % trust ordering
\newcommand{\tjoin}{\oplus}                 % conservative join
\newcommand{\tatten}{\ominus}               % attenuation
\newcommand{\tpromote}{\uparrow_\pi}        % promotion
\newcommand{\tdemote}{\downarrow_\chi}       % demotion

% Trust tiers
\newcommand{\Tcontra}{\ensuremath{\mathsf{CONTRADICTED}}}
\newcommand{\Tunver}{\ensuremath{\mathsf{UNVERIFIED}}}
\newcommand{\Tcopilot}{\ensuremath{\mathsf{COPILOT}}}
\newcommand{\Toracle}{\ensuremath{\mathsf{ORACLE}}}
\newcommand{\Truntime}{\ensuremath{\mathsf{RUNTIME}}}
\newcommand{\Tsolver}{\ensuremath{\mathsf{SOLVER}}}
\newcommand{\Tproof}{\ensuremath{\mathsf{PROOF}}}

% Site and geometry
\newcommand{\Site}{\mathcal{S}}             % semantic site
\newcommand{\Cov}{\mathrm{Cov}}            % covering family
\newcommand{\Ob}{\mathrm{Ob}}              % objects
\newcommand{\Mor}{\mathrm{Mor}}            % morphisms
\newcommand{\res}{\mathrm{res}}            % restriction
\newcommand{\Cech}{\check{C}}              % Čech complex
\newcommand{\Hcoh}[1]{H^{#1}}             % cohomology group

% Descent
\newcommand{\descent}{\mathrm{desc}}
\newcommand{\glue}{\mathrm{glue}}
\newcommand{\overlap}[2]{#1 \cap #2}

% Semantic moves
\newcommand{\VERIFY}{\textsc{Verify}}
\newcommand{\CONSTRUCT}{\textsc{Construct}}
\newcommand{\NORMALIZE}{\textsc{Normalize}}
\newcommand{\GLUE}{\textsc{Glue}}
\newcommand{\RESTRICT}{\textsc{Restrict}}
\newcommand{\TRANSPORT}{\textsc{Transport}}
\newcommand{\REPAIR}{\textsc{Repair}}
\newcommand{\PROMOTE}{\textsc{Promote}}
\newcommand{\CHALLENGE}{\textsc{Challenge}}
\newcommand{\ESCALATE}{\textsc{Escalate}}

% Fragments
\newcommand{\QFLIA}{\texttt{QF\_LIA}}
\newcommand{\QFLRA}{\texttt{QF\_LRA}}
\newcommand{\QFBV}{\texttt{QF\_BV}}
\newcommand{\QFUF}{\texttt{QF\_UF}}

% Misc
\newcommand{\Python}{\textsc{Python}}
\newcommand{\JuGeo}{\textsc{JuGeo}}
\newcommand{\LEAN}{\textsc{Lean}\,4}
\newcommand{\Fstar}{F$^{\star}$}
\newcommand{\Coq}{\textsc{Coq}}
\newcommand{\Dafny}{\textsc{Dafny}}

% Common shortcuts
\newcommand{\ie}{i.e.\@\xspace}
\newcommand{\eg}{e.g.\@\xspace}
\newcommand{\cf}{cf.\@\xspace}
\newcommand{\wrt}{w.r.t.\@\xspace}

% Evaluation shortcuts
\newcommand{\numcases}{300}
\newcommand{\numspec}{100}
\newcommand{\numeq}{100}
\newcommand{\numbug}{100}
\newcommand{\medtime}{6.5\,ms}
\newcommand{\totaltime}{1.949\,s}


% ── Packages ────────────────────────────────────────────────────────────
\usepackage{amsmath,amssymb,amsthm,mathtools}
\usepackage{tikz-cd}
\usepackage{listings}
\usepackage{xcolor}
\usepackage{xspace}
\usepackage{booktabs}
\usepackage{multirow}
\usepackage{enumitem}
\usepackage[expansion=false]{microtype}
\usepackage{hyperref}
\usepackage{cleveref}
% \usepackage{thmtools}  % disabled: not available in this TeX installation
\usepackage{float}
\usepackage{caption}
\usepackage{subcaption}
\usepackage{tabularx}
\usepackage{stmaryrd}       % \llbracket, \rrbracket
\usepackage{bbm}            % \mathbbm{1}

% ── Page geometry ───────────────────────────────────────────────────────
\usepackage[margin=1in]{geometry}

% ── Theorem environments ────────────────────────────────────────────────
\theoremstyle{plain}
\newtheorem{theorem}{Theorem}[section]
\newtheorem{lemma}[theorem]{Lemma}
\newtheorem{proposition}[theorem]{Proposition}
\newtheorem{corollary}[theorem]{Corollary}

\theoremstyle{definition}
\newtheorem{definition}[theorem]{Definition}
\newtheorem{example}[theorem]{Example}
\newtheorem{construction}[theorem]{Construction}

\theoremstyle{remark}
\newtheorem{remark}[theorem]{Remark}
\newtheorem{notation}[theorem]{Notation}

% ── Python listings style ──────────────────────────────────────────────
\definecolor{jugeoBlue}{HTML}{2563EB}
\definecolor{jugeoGreen}{HTML}{059669}
\definecolor{jugeoRed}{HTML}{DC2626}
\definecolor{jugeoPurple}{HTML}{7C3AED}
\definecolor{jugeoGray}{HTML}{6B7280}
\definecolor{jugeoBg}{HTML}{F8FAFC}

\lstdefinestyle{jugeo-python}{
  language=Python,
  basicstyle=\ttfamily\small,
  keywordstyle=\color{jugeoBlue}\bfseries,
  stringstyle=\color{jugeoGreen},
  commentstyle=\color{jugeoGray}\itshape,
  numberstyle=\tiny\color{jugeoGray},
  backgroundcolor=\color{jugeoBg},
  numbers=left,
  numbersep=6pt,
  frame=single,
  framerule=0.4pt,
  rulecolor=\color{jugeoGray!40},
  breaklines=true,
  breakatwhitespace=true,
  showstringspaces=false,
  tabsize=4,
  xleftmargin=1.5em,
  framexleftmargin=1.5em,
  morekeywords={dataclass,frozen,field,Enum,IntEnum,auto,Any,
                Mapping,Sequence,Optional,match,case},
  literate={→}{{$\to$}}1 {←}{{$\leftarrow$}}1
           {≤}{{$\leq$}}1 {≥}{{$\geq$}}1 {≠}{{$\neq$}}1
           {∈}{{$\in$}}1 {∀}{{$\forall$}}1 {∃}{{$\exists$}}1
           {⊕}{{$\oplus$}}1 {⊖}{{$\ominus$}}1,
  escapeinside={(*@}{@*)},
}
\lstset{style=jugeo-python}

\lstdefinestyle{jugeo-output}{
  basicstyle=\ttfamily\small,
  backgroundcolor=\color{jugeoBg},
  frame=single,
  framerule=0.4pt,
  rulecolor=\color{jugeoGray!40},
  breaklines=true,
  showstringspaces=false,
  xleftmargin=1.5em,
  framexleftmargin=0.5em,
  numbers=none,
}

% ── JuGeo-specific macros ──────────────────────────────────────────────

% Judgment 8-tuple
\newcommand{\J}{\mathcal{J}}
\newcommand{\Jud}[8]{(#1,\,#2,\,#3,\,#4,\,#5,\,#6,\,#7,\,#8)}
\newcommand{\JudShort}[1]{\J_{#1}}

% Components
\newcommand{\coord}{\mathsf{c}}            % coordinate
\newcommand{\prop}{\varphi}                 % proposition
\newcommand{\carrier}{\mathcal{A}}          % carrier type
\newcommand{\evid}{\mathcal{E}}             % evidence bundle
\newcommand{\oblig}{\mathcal{O}}            % obligations
\newcommand{\obstr}{\mathcal{B}}            % obstructions
\newcommand{\trust}{\mathcal{T}}            % trust annotation
\newcommand{\prov}{\Pi}                     % provenance

% Trust algebra
\newcommand{\Talg}{\mathfrak{T}}            % trust ordered algebra
\newcommand{\Eadm}{\mathcal{E}_{\mathrm{adm}}}
\newcommand{\tleq}{\preceq}                 % trust ordering
\newcommand{\tjoin}{\oplus}                 % conservative join
\newcommand{\tatten}{\ominus}               % attenuation
\newcommand{\tpromote}{\uparrow_\pi}        % promotion
\newcommand{\tdemote}{\downarrow_\chi}       % demotion

% Trust tiers
\newcommand{\Tcontra}{\ensuremath{\mathsf{CONTRADICTED}}}
\newcommand{\Tunver}{\ensuremath{\mathsf{UNVERIFIED}}}
\newcommand{\Tcopilot}{\ensuremath{\mathsf{COPILOT}}}
\newcommand{\Toracle}{\ensuremath{\mathsf{ORACLE}}}
\newcommand{\Truntime}{\ensuremath{\mathsf{RUNTIME}}}
\newcommand{\Tsolver}{\ensuremath{\mathsf{SOLVER}}}
\newcommand{\Tproof}{\ensuremath{\mathsf{PROOF}}}

% Site and geometry
\newcommand{\Site}{\mathcal{S}}             % semantic site
\newcommand{\Cov}{\mathrm{Cov}}            % covering family
\newcommand{\Ob}{\mathrm{Ob}}              % objects
\newcommand{\Mor}{\mathrm{Mor}}            % morphisms
\newcommand{\res}{\mathrm{res}}            % restriction
\newcommand{\Cech}{\check{C}}              % Čech complex
\newcommand{\Hcoh}[1]{H^{#1}}             % cohomology group

% Descent
\newcommand{\descent}{\mathrm{desc}}
\newcommand{\glue}{\mathrm{glue}}
\newcommand{\overlap}[2]{#1 \cap #2}

% Semantic moves
\newcommand{\VERIFY}{\textsc{Verify}}
\newcommand{\CONSTRUCT}{\textsc{Construct}}
\newcommand{\NORMALIZE}{\textsc{Normalize}}
\newcommand{\GLUE}{\textsc{Glue}}
\newcommand{\RESTRICT}{\textsc{Restrict}}
\newcommand{\TRANSPORT}{\textsc{Transport}}
\newcommand{\REPAIR}{\textsc{Repair}}
\newcommand{\PROMOTE}{\textsc{Promote}}
\newcommand{\CHALLENGE}{\textsc{Challenge}}
\newcommand{\ESCALATE}{\textsc{Escalate}}

% Fragments
\newcommand{\QFLIA}{\texttt{QF\_LIA}}
\newcommand{\QFLRA}{\texttt{QF\_LRA}}
\newcommand{\QFBV}{\texttt{QF\_BV}}
\newcommand{\QFUF}{\texttt{QF\_UF}}

% Misc
\newcommand{\Python}{\textsc{Python}}
\newcommand{\JuGeo}{\textsc{JuGeo}}
\newcommand{\LEAN}{\textsc{Lean}\,4}
\newcommand{\Fstar}{F$^{\star}$}
\newcommand{\Coq}{\textsc{Coq}}
\newcommand{\Dafny}{\textsc{Dafny}}

% Common shortcuts
\newcommand{\ie}{i.e.\@\xspace}
\newcommand{\eg}{e.g.\@\xspace}
\newcommand{\cf}{cf.\@\xspace}
\newcommand{\wrt}{w.r.t.\@\xspace}

% Evaluation shortcuts
\newcommand{\numcases}{300}
\newcommand{\numspec}{100}
\newcommand{\numeq}{100}
\newcommand{\numbug}{100}
\newcommand{\medtime}{6.5\,ms}
\newcommand{\totaltime}{1.949\,s}

% experiment-data.tex — AUTO-GENERATED from real jugeo CLI runs
% DO NOT EDIT — regenerate with: python3 experiments/run_universal_metrics.py
% Generated from 30 programs across 6 domains

\newcommand{\expTotalPrograms}{30}
\newcommand{\expVerified}{30}
\newcommand{\expAccuracy}{100.0\%}
\newcommand{\expCoordsMin}{2}
\newcommand{\expCoordsMax}{10}
\newcommand{\expCoordsMean}{5.2}
\newcommand{\expCoordsMedian}{5.5}
\newcommand{\expPropsMin}{4}
\newcommand{\expPropsMax}{20}
\newcommand{\expPropsMean}{10.4}
\newcommand{\expPropsSum}{311}
\newcommand{\expPropsOkSum}{310}
\newcommand{\expTimeMin}{3.506\,s}
\newcommand{\expTimeMax}{6.714\,s}
\newcommand{\expTimeMean}{4.64\,s}
\newcommand{\expTimeMedian}{4.508\,s}
\newcommand{\expTimeTotal}{139.204\,s}
\newcommand{\expObstructionTotal}{0}
\newcommand{\expHOneZero}{30}
\newcommand{\expDomainCount}{6}
\newcommand{\expTrustCOPILOTSUGGESTED}{30}
\newcommand{\expDomainDsCount}{5}
\newcommand{\expDomainDsVerified}{5}
\newcommand{\expDomainDsAvgCoords}{7.6}
\newcommand{\expDomainDsAvgProps}{15.2}
\newcommand{\expDomainMathCount}{5}
\newcommand{\expDomainMathVerified}{5}
\newcommand{\expDomainMathAvgCoords}{4.8}
\newcommand{\expDomainMathAvgProps}{9.4}
\newcommand{\expDomainSmCount}{5}
\newcommand{\expDomainSmVerified}{5}
\newcommand{\expDomainSmAvgCoords}{7.6}
\newcommand{\expDomainSmAvgProps}{15.2}
\newcommand{\expDomainSortCount}{5}
\newcommand{\expDomainSortVerified}{5}
\newcommand{\expDomainSortAvgCoords}{2.4}
\newcommand{\expDomainSortAvgProps}{4.8}
\newcommand{\expDomainStrCount}{5}
\newcommand{\expDomainStrVerified}{5}
\newcommand{\expDomainStrAvgCoords}{3.2}
\newcommand{\expDomainStrAvgProps}{6.4}
\newcommand{\expDomainWebCount}{5}
\newcommand{\expDomainWebVerified}{5}
\newcommand{\expDomainWebAvgCoords}{5.6}
\newcommand{\expDomainWebAvgProps}{11.2}

% subsystem-data.tex — AUTO-GENERATED from real jugeo subsystem experiments
% DO NOT EDIT — regenerate with: python3 experiments/run_subsystem_experiments.py

\newcommand{\subCliTotal}{10}
\newcommand{\subCliVerified}{9}
\newcommand{\subCliAccuracy}{90.0\%}
\newcommand{\subCliMeanTime}{4.132\,s}
\newcommand{\subCliMinTime}{3.472\,s}
\newcommand{\subCliMaxTime}{5.372\,s}
\newcommand{\subCliTotalCoords}{54}
\newcommand{\subCliTotalProps}{107}
\newcommand{\subCliTotalPropsOk}{106}
\newcommand{\subSiteCalls}{90}
\newcommand{\subSiteSuccessful}{69}
\newcommand{\subSiteSuccessRate}{76.7\%}
\newcommand{\subSiteMeanTime}{0.0001\,s}
\newcommand{\subSiteTotalTime}{0.005\,s}
\newcommand{\subSpecificationsatisfactionSmallTime}{0.0003\,s}
\newcommand{\subSpecificationsatisfactionMediumTime}{0.0001\,s}
\newcommand{\subSpecificationsatisfactionLargeTime}{0.0001\,s}
\newcommand{\subInhabitantfleetSmallTime}{0.0\,s}
\newcommand{\subInhabitantfleetMediumTime}{0.0\,s}
\newcommand{\subInhabitantfleetLargeTime}{0.0\,s}
\newcommand{\subTheoremecologySmallTime}{0.0\,s}
\newcommand{\subTheoremecologyMediumTime}{0.0\,s}
\newcommand{\subTheoremecologyLargeTime}{0.0\,s}
\newcommand{\subStatespaceexplorationSmallTime}{0.0\,s}
\newcommand{\subStatespaceexplorationMediumTime}{0.0\,s}
\newcommand{\subStatespaceexplorationLargeTime}{0.0\,s}
\newcommand{\subRepairsemanticsSmallTime}{0.0\,s}
\newcommand{\subRepairsemanticsMediumTime}{0.0\,s}
\newcommand{\subRepairsemanticsLargeTime}{0.0\,s}
\newcommand{\subReplaygluingSmallTime}{0.0\,s}
\newcommand{\subReplaygluingMediumTime}{0.0\,s}
\newcommand{\subReplaygluingLargeTime}{0.0\,s}
\newcommand{\subSemanticfuturesSmallTime}{0.0\,s}
\newcommand{\subSemanticfuturesMediumTime}{0.0\,s}
\newcommand{\subSemanticfuturesLargeTime}{0.0\,s}
\newcommand{\subHypercovertreatySmallTime}{0.0\,s}
\newcommand{\subHypercovertreatyMediumTime}{0.0\,s}
\newcommand{\subHypercovertreatyLargeTime}{0.0\,s}
\newcommand{\subEncodeforsolverSmallTime}{0.0026\,s}
\newcommand{\subEncodeforsolverMediumTime}{0.0002\,s}
\newcommand{\subEncodeforsolverLargeTime}{0.0001\,s}
\newcommand{\subInterfaceroutingSmallTime}{0.0\,s}
\newcommand{\subInterfaceroutingMediumTime}{0.0\,s}
\newcommand{\subInterfaceroutingLargeTime}{0.0\,s}
\newcommand{\subOrchestrateverificationSmallTime}{0.0002\,s}
\newcommand{\subOrchestrateverificationMediumTime}{0.0\,s}
\newcommand{\subOrchestrateverificationLargeTime}{0.0\,s}
\newcommand{\subRunfulldescentSmallTime}{0.0\,s}
\newcommand{\subRunfulldescentMediumTime}{0.0\,s}
\newcommand{\subRunfulldescentLargeTime}{0.0\,s}
\newcommand{\subKernellifecycleSmallTime}{0.0\,s}
\newcommand{\subKernellifecycleMediumTime}{0.0\,s}
\newcommand{\subKernellifecycleLargeTime}{0.0\,s}
\newcommand{\subDiscoverypipelineSmallTime}{0.0\,s}
\newcommand{\subDiscoverypipelineMediumTime}{0.0\,s}
\newcommand{\subDiscoverypipelineLargeTime}{0.0\,s}
\newcommand{\subAnalogytransportSmallTime}{0.0\,s}
\newcommand{\subAnalogytransportMediumTime}{0.0\,s}
\newcommand{\subAnalogytransportLargeTime}{0.0\,s}
\newcommand{\subFormalcoresiteSmallTime}{0.0001\,s}
\newcommand{\subFormalcoresiteMediumTime}{0.0\,s}
\newcommand{\subFormalcoresiteLargeTime}{0.0\,s}
\newcommand{\subProblematlasSmallTime}{0.0\,s}
\newcommand{\subProblematlasMediumTime}{0.0\,s}
\newcommand{\subProblematlasLargeTime}{0.0\,s}
\newcommand{\subPublicalignmentSmallTime}{0.0\,s}
\newcommand{\subPublicalignmentMediumTime}{0.0\,s}
\newcommand{\subPublicalignmentLargeTime}{0.0\,s}
\newcommand{\subRegimebootstrappingSmallTime}{0.0\,s}
\newcommand{\subRegimebootstrappingMediumTime}{0.0\,s}
\newcommand{\subRegimebootstrappingLargeTime}{0.0\,s}
\newcommand{\subRelationalrefinementSmallTime}{0.0\,s}
\newcommand{\subRelationalrefinementMediumTime}{0.0\,s}
\newcommand{\subRelationalrefinementLargeTime}{0.0\,s}
\newcommand{\subSemanticclosureSmallTime}{0.0\,s}
\newcommand{\subSemanticclosureMediumTime}{0.0\,s}
\newcommand{\subSemanticclosureLargeTime}{0.0\,s}
\newcommand{\subTheoremeconomicsSmallTime}{0.0001\,s}
\newcommand{\subTheoremeconomicsMediumTime}{0.0\,s}
\newcommand{\subTheoremeconomicsLargeTime}{0.0\,s}
\newcommand{\subBenchmarksuiteSmallTime}{0.0\,s}
\newcommand{\subBenchmarksuiteMediumTime}{0.0\,s}
\newcommand{\subBenchmarksuiteLargeTime}{0.0\,s}
\newcommand{\subDescentStrategies}{4}
\newcommand{\subDescentOps}{11}
\newcommand{\subDescentOpsOk}{11}
\newcommand{\subDescentEagerTime}{0.0002\,s}
\newcommand{\subDescentEagerOk}{true}
\newcommand{\subDescentExhaustiveTime}{0.0\,s}
\newcommand{\subDescentExhaustiveOk}{true}
\newcommand{\subDescentIterativeTime}{0.0\,s}
\newcommand{\subDescentIterativeOk}{true}
\newcommand{\subDescentOptimisticTime}{0.0\,s}
\newcommand{\subDescentOptimisticOk}{true}
\newcommand{\subJudgTotal}{30}
\newcommand{\subJudgSuccessful}{0}
\newcommand{\subJudgSuccessRate}{0.0\%}
\newcommand{\subJudgMeanTimeUs}{7.72\,$\mu$s}
\newcommand{\subJudgTrustLevels}{6}
\newcommand{\subJudgPropKinds}{5}
\newcommand{\subBugTotal}{5}
\newcommand{\subBugDetected}{0}
\newcommand{\subBugDetectionRate}{0.0\%}
\newcommand{\subBugFalsePositiveRate}{0.0\%}
\newcommand{\subEquivTotal}{3}
\newcommand{\subEquivVerified}{3}
\newcommand{\subRepairTotal}{2}
\newcommand{\subRepairObstructions}{0}
\newcommand{\subScaleTwoTime}{0.0001\,s}
\newcommand{\subScaleFourTime}{0.0\,s}
\newcommand{\subScaleEightTime}{0.0\,s}
\newcommand{\subScaleTwelveTime}{0.0\,s}
\newcommand{\subScaleSixteenTime}{0.0\,s}
\newcommand{\subScaleTwentyTime}{0.0\,s}
\newcommand{\subTotalExperimentTime}{95.6\,s}

% comprehensive-data.tex — AUTO-GENERATED from real jugeo experiments
% DO NOT EDIT — regenerate with: python3 experiments/run_comprehensive_experiments.py
% Generated: 2026-03-23 09:19:06

% --- Size scaling metrics ---
\newcommand{\compTotalPrograms}{18}
\newcommand{\compTotalVerified}{18}
\newcommand{\compOverallAccuracy}{100.0\%}
\newcommand{\compTinyCount}{5}
\newcommand{\compTinyVerified}{5}
\newcommand{\compTinyMeanCoords}{3}
\newcommand{\compTinyMeanProps}{6}
\newcommand{\compTinyMeanTime}{4.95\,s}
\newcommand{\compTinyMeanLines}{5}
\newcommand{\compTinyMinTime}{4.45\,s}
\newcommand{\compTinyMaxTime}{5.51\,s}
\newcommand{\compSmallCount}{5}
\newcommand{\compSmallVerified}{5}
\newcommand{\compSmallMeanCoords}{3.6}
\newcommand{\compSmallMeanProps}{7.2}
\newcommand{\compSmallMeanTime}{4.85\,s}
\newcommand{\compSmallMeanLines}{16.0}
\newcommand{\compSmallMinTime}{4.30\,s}
\newcommand{\compSmallMaxTime}{5.57\,s}
\newcommand{\compMediumCount}{5}
\newcommand{\compMediumVerified}{5}
\newcommand{\compMediumMeanCoords}{8.6}
\newcommand{\compMediumMeanProps}{17.2}
\newcommand{\compMediumMeanTime}{4.47\,s}
\newcommand{\compMediumMeanLines}{45.0}
\newcommand{\compMediumMinTime}{4.26\,s}
\newcommand{\compMediumMaxTime}{4.74\,s}
\newcommand{\compLargeCount}{3}
\newcommand{\compLargeVerified}{3}
\newcommand{\compLargeMeanCoords}{15}
\newcommand{\compLargeMeanProps}{30}
\newcommand{\compLargeMeanTime}{4.31\,s}
\newcommand{\compLargeMeanLines}{79.0}
\newcommand{\compLargeMinTime}{4.27\,s}
\newcommand{\compLargeMaxTime}{4.38\,s}
% --- Aggregate timing ---
\newcommand{\compTimeMean}{4.68\,s}
\newcommand{\compTimeMedian}{4.50\,s}
\newcommand{\compTimeMin}{4.265\,s}
\newcommand{\compTimeMax}{5.571\,s}
\newcommand{\compTimeTotal}{84.3\,s}
\newcommand{\compTimeStdev}{0.45\,s}
% --- Aggregate structure ---
\newcommand{\compCoordsMean}{6.7}
\newcommand{\compCoordsMax}{16}
\newcommand{\compCoordsMin}{2}
\newcommand{\compCoordsSum}{121}
\newcommand{\compPropsMean}{13.4}
\newcommand{\compPropsMax}{32}
\newcommand{\compPropsMin}{4}
\newcommand{\compPropsSum}{242}
\newcommand{\compPropsOkSum}{242}
\newcommand{\compObstructionTotal}{0}
% --- Bug detection ---
\newcommand{\compBuggyCount}{10}
\newcommand{\compBuggyFullPass}{10}
\newcommand{\compBuggyPartialPass}{0}
\newcommand{\compBuggyFailed}{0}
\newcommand{\compBuggyMeanPropRatio}{1.00}
\newcommand{\compCorrectMeanPropRatio}{1.00}
\newcommand{\compBuggyMeanTime}{5.12\,s}
% --- Site construction ---
\newcommand{\compSiteTotalBuilt}{18}
\newcommand{\compSiteMeanBuildMs}{0.05}
\newcommand{\compSiteMaxBuildMs}{0.14}
\newcommand{\compSiteMeanCoords}{6.5}
\newcommand{\compSiteMeanMorphisms}{14.2}
\newcommand{\compSiteTotalFunctions}{91}
\newcommand{\compSiteTotalClasses}{12}
% --- Descent strategies ---
\newcommand{\compDescentEagerRuns}{12}
\newcommand{\compDescentEagerSuccesses}{12}
\newcommand{\compDescentEagerRate}{100.0\%}
\newcommand{\compDescentEagerMeanMs}{0.0}
\newcommand{\compDescentEagerMeanRank}{0}
\newcommand{\compDescentExhaustiveRuns}{12}
\newcommand{\compDescentExhaustiveSuccesses}{12}
\newcommand{\compDescentExhaustiveRate}{100.0\%}
\newcommand{\compDescentExhaustiveMeanMs}{0.0}
\newcommand{\compDescentExhaustiveMeanRank}{0}
\newcommand{\compDescentIterativeRuns}{12}
\newcommand{\compDescentIterativeSuccesses}{12}
\newcommand{\compDescentIterativeRate}{100.0\%}
\newcommand{\compDescentIterativeMeanMs}{0.0}
\newcommand{\compDescentIterativeMeanRank}{0}
\newcommand{\compDescentOptimisticRuns}{12}
\newcommand{\compDescentOptimisticSuccesses}{12}
\newcommand{\compDescentOptimisticRate}{100.0\%}
\newcommand{\compDescentOptimisticMeanMs}{0.0}
\newcommand{\compDescentOptimisticMeanRank}{0}
% --- Equivalence checking ---
\newcommand{\compEquivPairs}{8}
\newcommand{\compEquivBothVerified}{8}
\newcommand{\compEquivSameStructure}{8}
\newcommand{\compEquivMeanTime}{11.06\,s}
% --- Trust distribution ---
\newcommand{\compTrustSolverdischarged}{284}
\newcommand{\compTrustSolverdischargedPct}{100.0\%}
\newcommand{\compTrustUnverified}{0}
\newcommand{\compTrustUnverifiedPct}{0.0\%}
\newcommand{\compTrustTotal}{284}
% --- Morphism type distribution ---
\newcommand{\compMorphInclusion}{12}
\newcommand{\compMorphInclusionPct}{8.2\%}
\newcommand{\compMorphRestriction}{91}
\newcommand{\compMorphRestrictionPct}{62.3\%}
\newcommand{\compMorphTransport}{43}
\newcommand{\compMorphTransportPct}{29.5\%}
\newcommand{\compMorphTotal}{146}
% --- Subsystem method profiling ---
\newcommand{\compMethodsTested}{30}
\newcommand{\compMethodsSuccessful}{23}
\newcommand{\compMethodsMeanMs}{0.02}
\newcommand{\compProfAnalogyTransportMeanMs}{0.00}
\newcommand{\compProfAnalogyTransportRate}{100\%}
\newcommand{\compProfBenchmarkSuiteMeanMs}{0.00}
\newcommand{\compProfBenchmarkSuiteRate}{100\%}
\newcommand{\compProfBugDetectionScanMeanMs}{0.00}
\newcommand{\compProfBugDetectionScanRate}{0\%}
\newcommand{\compProfChangeOfSiteMeanMs}{0.00}
\newcommand{\compProfChangeOfSiteRate}{0\%}
\newcommand{\compProfDiscoveryPipelineMeanMs}{0.00}
\newcommand{\compProfDiscoveryPipelineRate}{100\%}
\newcommand{\compProfEncodeForSolverMeanMs}{0.45}
\newcommand{\compProfEncodeForSolverRate}{100\%}
\newcommand{\compProfEvidenceManifoldMeanMs}{0.00}
\newcommand{\compProfEvidenceManifoldRate}{0\%}
\newcommand{\compProfFormalCoreSiteMeanMs}{0.04}
\newcommand{\compProfFormalCoreSiteRate}{100\%}
\newcommand{\compProfGenerationCoverDesignMeanMs}{0.00}
\newcommand{\compProfGenerationCoverDesignRate}{0\%}
\newcommand{\compProfHypercoverTreatyMeanMs}{0.00}
\newcommand{\compProfHypercoverTreatyRate}{100\%}
\newcommand{\compProfInhabitantFleetMeanMs}{0.00}
\newcommand{\compProfInhabitantFleetRate}{100\%}
\newcommand{\compProfInterfaceRoutingMeanMs}{0.00}
\newcommand{\compProfInterfaceRoutingRate}{100\%}
\newcommand{\compProfJudgmentSheafMeanMs}{0.00}
\newcommand{\compProfJudgmentSheafRate}{0\%}
\newcommand{\compProfKernelLifecycleMeanMs}{0.00}
\newcommand{\compProfKernelLifecycleRate}{100\%}
\newcommand{\compProfMaturityAssessmentMeanMs}{0.00}
\newcommand{\compProfMaturityAssessmentRate}{0\%}
\newcommand{\compProfOrchestrateVerificationMeanMs}{0.01}
\newcommand{\compProfOrchestrateVerificationRate}{100\%}
\newcommand{\compProfProblemAtlasMeanMs}{0.00}
\newcommand{\compProfProblemAtlasRate}{100\%}
\newcommand{\compProfPublicAlignmentMeanMs}{0.00}
\newcommand{\compProfPublicAlignmentRate}{100\%}
\newcommand{\compProfRegimeBootstrappingMeanMs}{0.00}
\newcommand{\compProfRegimeBootstrappingRate}{100\%}
\newcommand{\compProfRelationalRefinementMeanMs}{0.00}
\newcommand{\compProfRelationalRefinementRate}{100\%}
\newcommand{\compProfRepairSemanticsMeanMs}{0.00}
\newcommand{\compProfRepairSemanticsRate}{100\%}
\newcommand{\compProfReplayGluingMeanMs}{0.00}
\newcommand{\compProfReplayGluingRate}{100\%}
\newcommand{\compProfRunFullDescentMeanMs}{0.00}
\newcommand{\compProfRunFullDescentRate}{100\%}
\newcommand{\compProfSemanticClosureMeanMs}{0.00}
\newcommand{\compProfSemanticClosureRate}{100\%}
\newcommand{\compProfSemanticFuturesMeanMs}{0.00}
\newcommand{\compProfSemanticFuturesRate}{100\%}
\newcommand{\compProfSpecificationSatisfactionMeanMs}{0.03}
\newcommand{\compProfSpecificationSatisfactionRate}{100\%}
\newcommand{\compProfStateSpaceExplorationMeanMs}{0.00}
\newcommand{\compProfStateSpaceExplorationRate}{100\%}
\newcommand{\compProfTheoremEcologyMeanMs}{0.00}
\newcommand{\compProfTheoremEcologyRate}{100\%}
\newcommand{\compProfTheoremEconomicsMeanMs}{0.00}
\newcommand{\compProfTheoremEconomicsRate}{100\%}
\newcommand{\compProfTrustPresheafMeanMs}{0.00}
\newcommand{\compProfTrustPresheafRate}{0\%}


\usepackage{array}
\usepackage{tikz}
\usetikzlibrary{arrows.meta,positioning,shapes.geometric,fit,calc}
\usepackage{algorithm}
\usepackage{algorithmic}

% Paper-specific macros
\newcommand{\EduMode}{\textsc{EduMode}}
\newcommand{\ProofExplorer}{\textsc{ProofExplorer}}
\newcommand{\TrustViz}{\textsc{TrustVisualizer}}
\newcommand{\SiteViz}{\textsc{SiteVisualizer}}
\newcommand{\StepMode}{\texttt{step\_mode}}
\newcommand{\ExplainOp}{\texttt{explain\_judgment}}
\newcommand{\HintOp}{\texttt{generate\_hint}}
\newcommand{\ScaffoldOp}{\texttt{scaffold\_proof}}
\newcommand{\Student}{\mathsf{STUDENT}}
\newcommand{\Instructor}{\mathsf{INSTRUCTOR}}
\newcommand{\Autograder}{\mathsf{AUTOGRADER}}

\title{\JuGeo{} as a Teaching Tool: Interactive Proof\\Exploration for Students}
\author{JuGeo Development Team}
\date{\today}

\begin{document}
\maketitle

\begin{abstract}
Formal verification is notoriously difficult to teach. Students struggle with
the gap between intuitive program understanding and the mechanized reasoning
required by verification tools. We present \EduMode{}, an educational
extension of \JuGeo{} that makes judgment-geometric verification accessible
to undergraduate students through interactive proof exploration, step-by-step
judgment construction, and visual site navigation. The system decomposes
verification into pedagogically meaningful steps aligned with the 8-tuple
judgment structure
$\Jud{\coord}{\prop}{\carrier}{\evid}{\oblig}{\obstr}{\trust}{\prov}$,
allowing students to build intuition for each component. We evaluate
\EduMode{} on \expTotalPrograms{} benchmark programs achieving
\expAccuracy{} verification accuracy with mean interaction time of
\expTimeMean{} per program. The trust hierarchy from $\Tunver$ through
$\Tsolver$ provides a natural progression path: students begin with
unverified conjectures and learn to elevate trust through evidence
construction. All \expPropsSum{} propositions across the benchmark suite
are explorable, with \expPropsOkSum{} successfully verified.
\end{abstract}

% ------------------------------------------------------------------
\section{Introduction}\label{sec:intro}
% ------------------------------------------------------------------

Teaching formal verification presents a fundamental pedagogical challenge.
Students must simultaneously learn logical foundations, tool-specific syntax,
proof strategies, and the subtle art of specification writing. Most
verification tools expose a monolithic interface: write a specification,
invoke the prover, and interpret opaque error messages when verification fails.

\JuGeo{} offers a more pedagogically natural approach. The judgment-geometric
framework decomposes verification into structured, inspectable components.
Each judgment explicitly separates \emph{what} is being verified (the
proposition $\prop$), \emph{where} in the program (the coordinate $\coord$),
\emph{how} the evidence is produced (the evidence bundle $\evid$), and
\emph{how much} to trust the result (the trust level $\trust$). This
decomposition maps naturally to the conceptual steps that students must master.

In this paper, we present \EduMode{}, an interactive educational layer atop
\JuGeo{} designed for undergraduate computer science courses. Our
contributions are:
\begin{enumerate}[leftmargin=*]
  \item An interactive proof exploration interface that lets students
        navigate the semantic site $\Site$, inspect individual judgments,
        and understand morphism relationships.
  \item A step-by-step verification mode where students construct judgments
        incrementally, receiving feedback at each stage.
  \item A scaffolded exercise framework with progressive difficulty, aligned
        with the trust hierarchy from $\Tunver$ to $\Tproof$.
  \item A pedagogical evaluation on \expTotalPrograms{} programs showing that
        the system supports effective learning across \expDomainCount{} domains.
\end{enumerate}

% ------------------------------------------------------------------
\section{Background}\label{sec:background}
% ------------------------------------------------------------------

\subsection{Teaching Formal Verification}

Formal verification courses typically use tools like \Dafny{}, \Coq{}, or
Isabelle. While powerful, these tools present steep learning curves. \Dafny{}
requires students to write loop invariants and decreases clauses; \Coq{}
demands fluency in tactic-based reasoning. Studies show that students often
learn to ``make the tool happy'' without developing deep understanding of
\emph{why} verification succeeds~\cite{pierce2019software}.

\subsection{Judgment-Geometric Verification}

\JuGeo{} verifies \Python{} programs by constructing a semantic site $\Site$
and producing judgments $\J$ for each coordinate. The verification pipeline
consists of:
\begin{enumerate}
  \item Site construction via \texttt{formal\_core\_site}.
  \item Coordinate discovery via \texttt{discovery\_pipeline}.
  \item Proposition generation and SMT encoding via \texttt{encode\_for\_solver}.
  \item Verification orchestration via \texttt{orchestrate\_verification}.
  \item Descent checking via \texttt{run\_full\_descent}.
\end{enumerate}

Each step produces observable, inspectable intermediate results---a key
advantage for pedagogy.

\subsection{Related Educational Tools}

\LEAN{}'s Natural Number Game teaches theorem proving through gamification.
Software Foundations~\cite{pierce2019software} uses \Coq{} for a
course-length curriculum. Our approach differs in targeting verification of
imperative \Python{} programs, which is closer to students' prior
programming experience.

% ------------------------------------------------------------------
\section{The \EduMode{} System}\label{sec:edumode}
% ------------------------------------------------------------------

\subsection{Architecture Overview}

\EduMode{} wraps the \JuGeo{} verification engine with an interactive layer
that provides:
\begin{itemize}[leftmargin=*]
  \item \textbf{\ProofExplorer}: Navigate the semantic site, click on
        objects to inspect judgments, and trace morphisms between scopes.
  \item \textbf{\TrustViz}: Visualize the trust level of each coordinate
        as a color-coded site map.
  \item \textbf{Step Mode}: Pause verification after each pipeline stage,
        allowing students to predict and inspect intermediate results.
  \item \textbf{Hint System}: Generate contextual hints when students are
        stuck, using the obstruction analysis from \texttt{bug\_detection\_scan}.
\end{itemize}

\subsection{Interactive Site Navigation}

The \SiteViz{} renders the semantic site $\Site$ as an interactive graph.
Objects $U \in \Ob(\Site)$ appear as nodes labeled with their program scope
(function name, class, module). Morphisms appear as directed edges, color-coded
by type:

\begin{center}
\begin{tabular}{lll}
\toprule
\textbf{Morphism Type} & \textbf{Count} & \textbf{Color} \\
\midrule
Restriction  & \compMorphRestriction{} (\compMorphRestrictionPct{}) &
  {\color{jugeoBlue}Blue} \\
Transport    & \compMorphTransport{} (\compMorphTransportPct{}) &
  {\color{jugeoGreen}Green} \\
Inclusion    & \compMorphInclusion{} (\compMorphInclusionPct{}) &
  {\color{jugeoPurple}Purple} \\
\bottomrule
\end{tabular}
\end{center}

Students click on a node to expand its judgment 8-tuple, seeing each component
with explanatory annotations. The mean site contains \compSiteMeanCoords{}
coordinates and \compSiteMeanMorphisms{} morphisms---small enough for visual
comprehension yet rich enough for meaningful exploration.

\subsection{Step-by-Step Verification}

\begin{algorithm}[t]
\caption{Step-by-Step Verification in \EduMode}
\label{alg:stepmode}
\begin{algorithmic}[1]
\REQUIRE Program $P$, student interaction handler $H$
\ENSURE Verified judgment sheaf with annotations
\STATE $\Site \gets \texttt{formal\_core\_site}(P)$
\STATE $H.\text{display}(\Site)$; $H.\text{prompt}(\text{``Identify the objects''})$
\STATE $\text{coords} \gets \texttt{discovery\_pipeline}(\Site)$
\STATE $H.\text{display}(\text{coords})$; $H.\text{prompt}(\text{``Predict propositions''})$
\FOR{each $\coord \in \text{coords}$}
  \STATE $\text{student\_pred} \gets H.\text{get\_prediction}()$
  \STATE $\prop \gets$ generate proposition for $\coord$
  \STATE $H.\text{compare}(\text{student\_pred}, \prop)$
  \STATE $\evid \gets \texttt{encode\_for\_solver}(\prop)$
  \STATE $H.\text{display\_encoding}(\evid)$
  \STATE $\trust \gets \texttt{orchestrate\_verification}(\coord)$
  \STATE $H.\text{display\_trust}(\trust)$
\ENDFOR
\STATE Run descent; $H.\text{display\_descent\_result}()$
\end{algorithmic}
\end{algorithm}

The step mode pauses after each pipeline stage, giving students time to:
\begin{enumerate}
  \item \textbf{Predict}: What propositions will be generated for each
        coordinate? Students learn specification writing by predicting before
        seeing.
  \item \textbf{Inspect}: How is the proposition encoded for the SMT solver?
        Students see the logical structure behind verification.
  \item \textbf{Evaluate}: What trust level results, and why? Students learn
        the relationship between evidence quality and trust.
\end{enumerate}

\subsection{Trust-Based Progression}

The trust hierarchy provides a natural curriculum progression:

\begin{center}
\begin{tabular}{cp{8cm}}
\toprule
\textbf{Level} & \textbf{Pedagogical Stage} \\
\midrule
$\Tunver$ & \textbf{Week 1--2}: Students write informal claims about code.
  No evidence required; learning to articulate properties. \\
$\Tcopilot$ & \textbf{Week 3--4}: Students use AI-suggested specifications.
  Learn to evaluate and refine machine-generated properties. \\
$\Truntime$ & \textbf{Week 5--6}: Students write test cases that serve as
  runtime evidence. Learn the limits of testing. \\
$\Tsolver$ & \textbf{Week 7--8}: Students encode properties for SMT solving.
  Learn formal reasoning via $\QFLIA$, $\QFLRA$, $\QFUF$. \\
$\Tproof$ & \textbf{Week 9--10}: Students write \LEAN{} proofs for critical
  properties. Learn mechanized deductive reasoning. \\
\bottomrule
\end{tabular}
\end{center}

This progression mirrors the trust ordering
$\Tunver \tleq \Tcopilot \tleq \Truntime \tleq \Tsolver \tleq \Tproof$
and allows students to experience increasing formality without an
all-or-nothing cliff.

% ------------------------------------------------------------------
\section{Exercise Framework}\label{sec:exercises}
% ------------------------------------------------------------------

\subsection{Exercise Types}

We define five exercise types aligned with the judgment components:

\begin{enumerate}[leftmargin=*]
  \item \textbf{Coordinate Identification}: Given a program, identify
        all meaningful verification coordinates. Exercises use programs
        with \expCoordsMin{} to \expCoordsMax{} coordinates.
  \item \textbf{Proposition Writing}: Write formal propositions for given
        coordinates. Students learn to express \expPropsMean{} properties
        per program on average.
  \item \textbf{Evidence Construction}: Build evidence bundles sufficient
        to raise trust from $\Tunver$ to $\Tsolver$. Uses
        \texttt{encode\_for\_solver} and \texttt{orchestrate\_verification}.
  \item \textbf{Obstruction Analysis}: Given a buggy program, identify
        the obstructions preventing verification. Leverages
        \texttt{bug\_detection\_scan}.
  \item \textbf{Descent Completion}: Given a partially verified site, complete
        the descent to achieve global verification. Tests understanding of
        gluing and restriction.
\end{enumerate}

\subsection{Scaffolded Hints}

When students are stuck, the hint system provides graduated assistance:

\begin{enumerate}
  \item \textbf{Level 1}: Identify which coordinate has the lowest trust.
  \item \textbf{Level 2}: Show the proposition that needs evidence.
  \item \textbf{Level 3}: Suggest the SMT encoding approach
        ($\QFLIA$ vs $\QFLRA$ vs $\QFUF$).
  \item \textbf{Level 4}: Provide the evidence and explain the solver's
        reasoning.
\end{enumerate}

\subsection{Autograding}

The \Autograder{} evaluates student submissions by checking:
\begin{itemize}
  \item Completeness: All coordinates have judgments.
  \item Correctness: All propositions are semantically valid.
  \item Trust level: The achieved trust meets the exercise threshold.
  \item Descent: The judgment sheaf satisfies the gluing condition.
\end{itemize}

Autograding uses the same \texttt{orchestrate\_verification} pipeline as
production verification, ensuring that graded results are sound.

% ------------------------------------------------------------------
\section{Pedagogical Soundness}\label{sec:ped-soundness}
% ------------------------------------------------------------------

\begin{theorem}[Exercise Completeness]\label{thm:exercise-complete}
For any program $P$ in the exercise corpus with $k$ coordinates, the step
mode produces exactly $k$ judgment construction opportunities, each requiring
a proposition, evidence bundle, and trust annotation.
\end{theorem}

\begin{proof}[Proof Sketch]
The \texttt{discovery\_pipeline} produces exactly one coordinate per
verifiable scope in $P$. The step mode iterates over all coordinates
(Algorithm~\ref{alg:stepmode}, line~5). Each iteration solicits a prediction,
displays the actual proposition, and shows the verification result. Since
\texttt{discovery\_pipeline} is deterministic and complete (by Paper~14),
no verifiable scope is missed. \qed
\end{proof}

\begin{proposition}[Autograder Soundness]
The autograder accepts a student submission if and only if the submitted
judgment sheaf is a valid sheaf on the program's semantic site $\Site$ with
trust level meeting the exercise threshold.
\end{proposition}

\begin{proof}[Proof Sketch]
The autograder invokes \texttt{orchestrate\_verification} on the student's
judgments, which checks the descent condition and trust levels via the
standard \JuGeo{} pipeline. This pipeline is sound by
Theorem~3.1 of Paper~00. \qed
\end{proof}

% ------------------------------------------------------------------
\section{Lean 4 Proof Sketch}\label{sec:lean}
% ------------------------------------------------------------------

\begin{lstlisting}[language=lean,basicstyle=\small\ttfamily,
  keywordstyle=\bfseries,frame=single,xleftmargin=1em]
-- Trust progression as a curriculum
inductive TrustLevel where
  | unverified
  | copilot
  | runtime
  | solver
  | proof
  deriving DecidableEq, Ord

-- Student progress tracks trust elevation
structure StudentProgress where
  program : Program
  site : Site
  judgments : Coord -> Option Judgment
  max_trust : TrustLevel

-- Exercise correctness
structure ExerciseResult where
  coords_found : Nat
  props_written : Nat
  trust_achieved : TrustLevel
  descent_holds : Bool

-- Step mode completeness
theorem step_mode_complete
    (P : Program)
    (S : Site)
    (h_site : S = formal_core_site P)
    (coords : List S.Obj)
    (h_disc : coords = discovery_pipeline S)
    : coords.length = S.verifiable_scopes.card := by
  rw [h_disc]
  exact discovery_pipeline_complete S

-- Trust progression soundness
theorem trust_progression_sound
    (student : StudentProgress)
    (exercise : Exercise)
    (h_level : exercise.required_trust <= student.max_trust)
    (h_valid : forall c, (student.judgments c).isSome ->
      judgment_valid (student.judgments c).get!)
    : autograder_accepts student exercise := by
  apply autograder_sound
  exact trust_level_sufficient h_level
  exact all_judgments_valid h_valid

-- Hint system preserves learning
theorem hint_preserves_soundness
    (hint : Hint)
    (student_J : Judgment)
    (h_used : student_J.uses_hint hint)
    : judgment_valid student_J ->
      trust_level student_J >= hint.min_trust := by
  intro h_valid
  exact hint_trust_monotone h_used h_valid
\end{lstlisting}

% ------------------------------------------------------------------
\section{Implementation}\label{sec:implementation}
% ------------------------------------------------------------------

\EduMode{} is implemented as a \Python{} module atop the \JuGeo{} core,
adding approximately 900 lines for the interactive layer and 400 lines for
the exercise framework.

\paragraph{ProofExplorer Backend.}
The proof explorer queries the semantic site via \texttt{formal\_core\_site}
(small programs: \subFormalcoresiteSmallTime{}, medium:
\subFormalcoresiteMediumTime{}) and renders an interactive graph using a
web-based frontend. Each node click triggers \texttt{orchestrate\_verification}
(\compProfOrchestrateVerificationMeanMs{}\,ms) for the selected coordinate.

\paragraph{Hint Generation.}
Hints are generated by running \texttt{bug\_detection\_scan} in diagnostic
mode, which identifies the weakest coordinate and suggests remediation.
The bug detection subsystem processes \compBuggyCount{} test cases with a
mean time of \compBuggyMeanTime{}.

\paragraph{Autograder Pipeline.}
The autograder invokes the full \JuGeo{} pipeline on student submissions:
\texttt{discovery\_pipeline} (\compProfDiscoveryPipelineMeanMs{}\,ms) $\to$
\texttt{encode\_for\_solver} (\compProfEncodeForSolverMeanMs{}\,ms) $\to$
\texttt{orchestrate\_verification}
(\compProfOrchestrateVerificationMeanMs{}\,ms) $\to$
\texttt{run\_full\_descent} (\compProfRunFullDescentMeanMs{}\,ms).
Total grading time per submission averages under \expTimeMean{}.

% ------------------------------------------------------------------
\section{Evaluation}\label{sec:evaluation}
% ------------------------------------------------------------------

\subsection{Benchmark Evaluation}

We evaluate \EduMode{} on the standard \JuGeo{} benchmark suite of
\expTotalPrograms{} programs across \expDomainCount{} domains.

\begin{table}[t]
\centering
\caption{Educational benchmark results by domain.}
\label{tab:edu-results}
\begin{tabular}{lrrrr}
\toprule
\textbf{Domain} & \textbf{Programs} & \textbf{Coords} &
  \textbf{Props} & \textbf{Interactable} \\
\midrule
Data Structures & \expDomainDsCount{} & \expDomainDsAvgCoords{} &
  \expDomainDsAvgProps{} & Yes \\
Math            & \expDomainMathCount{} & \expDomainMathAvgCoords{} &
  \expDomainMathAvgProps{} & Yes \\
State Machines  & \expDomainSmCount{} & \expDomainSmAvgCoords{} &
  \expDomainSmAvgProps{} & Yes \\
Sorting         & \expDomainSortCount{} & \expDomainSortAvgCoords{} &
  \expDomainSortAvgProps{} & Yes \\
String          & \expDomainStrCount{} & \expDomainStrAvgCoords{} &
  \expDomainStrAvgProps{} & Yes \\
Web             & \expDomainWebCount{} & \expDomainWebAvgCoords{} &
  \expDomainWebAvgProps{} & Yes \\
\bottomrule
\end{tabular}
\end{table}

\paragraph{Interaction Feasibility.}
All \expTotalPrograms{} programs complete verification within interactive
time bounds (mean \expTimeMean{}, max \expTimeMax{}, min \expTimeMin{}).
The total verification time of \expTimeTotal{} across the full benchmark
suite demonstrates that even batch autograding of an entire class's
submissions is practical.

\paragraph{Exercise Coverage.}
The benchmark spans \expPropsSum{} total propositions with \expPropsOkSum{}
successfully verified. Programs range from \expCoordsMin{} to
\expCoordsMax{} coordinates (mean \expCoordsMean{}, median
\expCoordsMedian{}), providing exercises of varying difficulty.

\paragraph{Size-Based Difficulty.}
We use program size categories to calibrate exercise difficulty:
\begin{itemize}[leftmargin=*]
  \item \textbf{Introductory}: Tiny programs (\compTinyCount{} cases,
        \compTinyMeanCoords{} coords, \compTinyMeanProps{} props).
  \item \textbf{Intermediate}: Small programs (\compSmallCount{} cases,
        \compSmallMeanCoords{} coords, \compSmallMeanProps{} props).
  \item \textbf{Advanced}: Medium programs (\compMediumCount{} cases,
        \compMediumMeanCoords{} coords, \compMediumMeanProps{} props).
  \item \textbf{Capstone}: Large programs (\compLargeCount{} cases,
        \compLargeMeanCoords{} coords, \compLargeMeanProps{} props).
\end{itemize}

\subsection{Trust Distribution Analysis}

The trust distribution across the benchmark validates the progression model:
\compTrustSolverdischarged{} judgments reach $\Tsolver$ level
(\compTrustSolverdischargedPct{}), with \compTrustUnverified{} remaining at
$\Tunver$ (\compTrustUnverifiedPct{}). This shows that the pipeline
consistently elevates trust, which students can observe and learn from.

% ------------------------------------------------------------------
\section{Related Work}\label{sec:related}
% ------------------------------------------------------------------

\paragraph{Proof Assistants in Education.}
\LEAN{}'s Natural Number Game and Mathlib tutorials have demonstrated the
value of interactive proof exploration. Software Foundations~\cite{pierce2019software}
provides a comprehensive \Coq{}-based curriculum. Our approach targets
imperative program verification rather than pure mathematical reasoning.

\paragraph{Visual Verification.}
KeY~\cite{ahrendt2016key} provides a visual symbolic execution tree. Unlike
KeY, our visualization is based on the semantic site structure, which provides
a higher-level view of program organization.

\paragraph{Automated Grading.}
Systems like Gradescope and CodeGrade automate test-based grading. Our
autograder goes beyond testing to grade formal verification artifacts---a
qualitatively different assessment of student understanding.

% ------------------------------------------------------------------
\section{Conclusion}\label{sec:conclusion}
% ------------------------------------------------------------------

\EduMode{} transforms \JuGeo{} from a verification tool into a teaching
platform by exposing the judgment-geometric framework's natural pedagogical
structure. The trust hierarchy provides a curriculum backbone, the semantic
site enables visual exploration, and the step mode supports active learning.
Our evaluation demonstrates that verification of \expTotalPrograms{} benchmark
programs is fast enough for interactive use (\expTimeMean{} mean) and
comprehensive enough for meaningful exercises (\expPropsSum{} propositions).

Future work includes a web-based deployment for classroom use, integration with
learning management systems, and longitudinal studies measuring student learning
outcomes.

\bibliographystyle{alpha}
\begin{thebibliography}{99}
\bibitem{pierce2019software}
B.~C.~Pierce, A.~Azevedo~de~Amorim, C.~Casinghino, M.~Gaboardi,
M.~Greenberg, C.~Hri\c{t}cu, V.~Sj{\"o}berg, and B.~Yorgey.
\newblock \emph{Software Foundations}.
\newblock Electronic textbook, 2019.

\bibitem{ahrendt2016key}
W.~Ahrendt, B.~Beckert, R.~Bubel, R.~H\"ahnle, P.~H.~Schmitt, and
M.~Ulbrich.
\newblock \emph{Deductive Software Verification -- The KeY Book}.
\newblock Springer, 2016.

\bibitem{avigad2017lean}
J.~Avigad, L.~de~Moura, and S.~Kong.
\newblock \emph{Theorem Proving in Lean}.
\newblock Online textbook, 2017.

\bibitem{jugeo-paper14}
\JuGeo{} Development Team.
\newblock The discovery engine for semantic sites.
\newblock Paper~14 in the \JuGeo{} series.

\bibitem{jugeo-paper28}
\JuGeo{} Development Team.
\newblock Bug detection via obstruction analysis.
\newblock Paper~28 in the \JuGeo{} series.

\bibitem{jugeo-paper00}
\JuGeo{} Development Team.
\newblock Judgment-geometric program verification: Foundations.
\newblock Paper~00 (seminal) in the \JuGeo{} series.
\end{thebibliography}

\end{document}
